\section{The invariant Penrose inequality in four spatial dimensions}
\label{four:sec:introduction}

We prove a lower bound for the ADM rest mass in terms of the least
three-volume needed to separate a chosen asymptotically flat end from
the compact weakly trapped boundary and the other ends. The second
fundamental form is arbitrary. The final comparison uses the
Riemannian Penrose inequality in dimension four after a graph and
conformal deformation; the required enclosure is constructed below.

\subsection{Initial data and enclosing cuts}

Four is the spatial dimension; the associated spacetime dimension is
five, and a hypersurface area is a three-dimensional Riemannian volume.
Set
\[
 \omega=\omega_3=|S^3|=2\pi^2,\qquad
 \tau=\tr_gK,\qquad
 \mu=\tfrac12(R_g+\tau^2-|K|_g^2),\qquad
 J_i=\nabla^j(K_{ij}-\tau g_{ij}).
\]
Here $R_g=\Scal_g$ denotes scalar curvature and $\nabla$ is the
Levi--Civita connection of $g$. The cosmological constant is zero. We use
\(\Delta_g=\divg_g\grad_g\).
For a hypersurface with specified unit normal \(\nu\), write
\[
 H=\divg_{\mathrm{tan}}\nu,\qquad
 \tr_{\mathrm{tan}}K=(g^{ij}-\nu^i\nu^j)K_{ij},\qquad
 \Theta_K=H+\tr_{\mathrm{tan}}K.
\]
At an inner boundary the normal points into the exterior. In an auxiliary
trapped-region construction it points out of the bounded region.

The notation \(O_j(r^{-a})\) includes coordinate derivatives of order
\(k\le j\), bounded by \(O(r^{-a-k})\). On an asymptotically flat end the
ADM normalization is
\begin{align}
 E&=\frac{1}{6\omega}\lim_{R\to\infty}
       \int_{|x|=R}(\partial_jg_{ij}-\partial_i g_{jj})
                    n_\delta^i\,\dd A_\delta,\label{four:intro:energy}\\
 P_i&=\frac{1}{3\omega}\lim_{R\to\infty}
       \int_{|x|=R}(K_{ij}-\tau g_{ij})
                    n_\delta^j\,\dd A_\delta.\label{four:intro:momentum}
\end{align}
Here \(n_\delta\) and \(\dd A_\delta\) are Euclidean.
Whenever a charge is used below, its stated limit is required to exist.

\begin{definition}[Enclosing cuts]\label{four:def:cuts}
Let \(M\) have nonempty compact boundary \(S\) and finitely many
asymptotically flat ends, and fix an end \(e\).
An admissible outer domain \(D_e\) is a connected smooth
codimension-zero submanifold with boundary, closed as a subset of \(M\),
such that its manifold interior lies in \(\operatorname{int}M\),
it contains the entire sufficiently distant part of \(e\), and
it contains no sufficiently distant part of any other end.
Its entire intrinsic manifold boundary
\(\Gamma=\partial_{\mathrm{man}}D_e\) is required to be compact.
The enclosing volume is
\[
 A_e(S)=\inf_{D_e}|\partial_{\mathrm{man}}D_e|_g.
\]
All frontier coincident with \(S\) is counted. The cut may be
disconnected. In the one-ended case \(D_e=M\) is allowed and has
intrinsic boundary \(S\). We then also write \(A_*\) for this infimum.
\end{definition}

Only these smooth cuts are used. Their infimum need not be attained.
The definition treats the entire original boundary and all other ends
as obstacles to the selected outer domain.

\begin{theorem}\label{four:thm:main}
Let \((M^4,g,K)\) be smooth and
connected, with \(M\) orientable, \(K\) an arbitrary smooth symmetric
covariant two-tensor, and \(S=\partial M\) nonempty, compact, and smooth.
Suppose \(g\) is complete as a metric space with \(S\) included.
Outside a compact set let \(M\) have finitely many, and at least one,
ends diffeomorphic to the complement of a closed ball in \(\R^4\).
Assume
\[
 \mu\ge|J|_g,\qquad \mu,|J|_g\in L^1(M,\dd V_g),
\]
and, on every end for a common \(q>1\),
\[
 g-\delta=O_2(r^{-q}),\qquad K=O_1(r^{-1-q}),
\]
with finite energy and momentum limits
\eqref{four:intro:energy}--\eqref{four:intro:momentum}.
On every component of \(S\) assume \(\Theta_K\le0\), with normal into \(M\).

Then for each end \(e\),
\begin{equation}\label{four:intro:main}
 E_e\ge
 \sqrt{|P_e|_\delta^2+
       \frac14\left(\frac{A_e(S)}{\omega}\right)^{4/3}}.
\end{equation}
In particular \(E_e>|P_e|_\delta\), and the invariant mass satisfies
\begin{equation}\label{four:intro:mass}
 \sqrt{E_e^2-|P_e|_\delta^2}
   \ge\frac12\left(\frac{A_e(S)}{\omega}\right)^{2/3}.
\end{equation}
\end{theorem}

The strict timelikeness is a conclusion: Lemma~\ref{four:lem:cut-reduction}
establishes \(A_e(S)>0\), and the proof includes the non-timelike case.
There is no extra decay assumption on \(\tau=\tr_gK\).
All original boundary components use the same future trapping convention.

The coefficient is sharp. On the time-symmetric
Schwarzschild--Tangherlini exterior \cite{Tangherlini1963},
\[
 g_m=\left(1+\frac{m}{2r^2}\right)^2\delta,\qquad
 K=0,\qquad r\ge\sqrt{m/2},
\]
the ADM energy is \(m\), the momentum is zero, and the boundary
has three-volume \(\omega(2m)^{3/2}\).
The areal radius \(r+m/(2r)\) is nondecreasing on this exterior.
Radial projection to its boundary is therefore nonexpanding on
three-dimensional tangent volumes; its degree on an enclosing cut is one.
The minimum enclosing volume is consequently the boundary volume,
and \eqref{four:intro:main} is an equality for this family.
The equality statement below concerns the original exterior data and has
additional horizon hypotheses; these are not assumptions of the numerical
inequality.

\subsection{Equality for a connected exterior}

We also write $\theta_+=\Theta_K$ for the future null expansion.

\begin{definition}[The connected outermost horizon class]\label{rig4:def:horizon}
Let \((\Omega^4,g,K)\) satisfy the initial-data hypotheses of
Theorem~\ref{four:thm:main}, with exactly one end
and connected boundary \(S\). We additionally require:
\begin{enumerate}[label=(\roman*),leftmargin=*]
\item \(S\) is a future marginally outer trapped surface (MOTS):
      \(\theta_+(S)=0\) for the normal into \(\Omega\).
\item \(S\) is outermost toward the end: no smooth compact embedded two-sided
      enclosing hypersurface in \(\operatorname{int}\Omega\), possibly
      disconnected and oriented toward the end, has \(\theta_+\le0\)
      everywhere.
\item \(S\) is outer area-minimizing: every cut in Definition~\ref{four:def:cuts}
      has area at least \(A:=|S|_g>0\). Consequently \(A_e(S)=A\).
\end{enumerate}
\end{definition}

For a geometric mass \(M_0>0\), the Schwarzschild--Tangherlini metric
\cite{Tangherlini1963} in its static exterior is
\begin{equation}\label{rig4:eq:intro-tangherlini}
 \overline g_{M_0}
 =-\left(1-\frac{2M_0}{r^2}\right)dt^2
  +\left(1-\frac{2M_0}{r^2}\right)^{-1}dr^2+r^2g_{S^3},
 \qquad r>\sqrt{2M_0},
\end{equation}
where \(g_{S^3}\) is the unit round metric. ``Regular maximal extension''
refers to the usual extension through the future and past horizons,
including their bifurcation sphere. The apparent singularity in
Equation~\eqref{rig4:eq:intro-tangherlini} at the horizon is a coordinate singularity.

\begin{theorem}[Original-data exterior rigidity]\label{rig4:thm:rigidity}
Let \((\Omega,g,K)\) belong to Definition~\ref{rig4:def:horizon}, and put
\(m=\sqrt{E^2-|P|_\delta^2}\). Then
\begin{equation}\label{rig4:eq:intro-equality}
 m=\frac12\left(\frac A\omega\right)^{2/3}
\end{equation}
if and only if the original initial data admit a global smooth spacelike
embedding, including \(S\), into the regular maximal extension of
\(\overline g_m\). The image lies in the closure of the corresponding exterior,
the chosen end approaches its spatial infinity, and \(S\) maps to a smooth
section of its future horizon or to its bifurcation sphere. The pullback of
the spacetime metric is \(g\); for a consistent future unit normal \(n\),
the second fundamental form is precisely the original tensor \(K\), with
the convention
\begin{equation}\label{rig4:eq:intro-second-fundamental-form}
 K(U,V)=\overline g_m(\overline\nabla_U n,V).
\end{equation}
Both the embedding and the normal extend smoothly to \(S\) in horizon-regular
coordinates.

Equivalently, every such Schwarzschild--Tangherlini hypersurface in
$\overline g_{M_0}$ satisfying
Definition~\ref{rig4:def:horizon} and the prescribed asymptotic hypotheses has
invariant ADM mass \(M_0\), horizon area
\(\omega(2M_0)^{3/2}\), and equality in
Theorem~\ref{four:thm:main}.
\end{theorem}

No spherical topology, simple connectivity, stationary development, or
vacuum condition is assumed in Theorem~\ref{rig4:thm:rigidity}. The conclusion
permits non-time-symmetric and asymptotically boosted hypersurfaces. It
asserts rigidity only of the exterior in Definition~\ref{rig4:def:horizon} and
makes no assertion about data behind \(S\). A graph in the singular static
time coordinate is not required at the horizon.

\subsection{Proof strategy and parameter order}

The proof first replaces a timelike asymptotic end by a rest end,
while preserving the dominant energy condition and comparing every
enclosing cut. Strictification then gives positive energy slack and
strict boundary trapping. Two threshold frontiers provide boundary
collars and a small interval for the prescribed graph trace.

On the resulting exterior we choose a positive smooth coefficient
function $l:\R\to(0,\infty)$, solve for real functions $f,Z$, define
$t$ from them, and construct
\[
 \widehat g=e^{2t}\bigl(g+l(t)^2\dd f\otimes\dd f\bigr).
\]
The scalar identity retains the unrestricted term $-F\tau$, where
$F$ is the prescribed scalar trace in the graph equation.
A polynomial penalty controls it in the first-floor estimate; the
later height bounds let strict DEC slack absorb it in the final scalar
comparison. The coupled existence proof uses a scalar principal matrix
with three transverse eigenvalues equal to one. Its measurable axial
coefficient is handled by the regularity theorem proved in
Section~\ref{four:sec:regularity}.

The deformation produces nonnegative scalar curvature, negative boundary
mean curvature, and a vanishing upper bound for the energy change.
A separate perimeter minimization supplies an outer-minimizing minimal
enclosure, so Bray--Lee's Riemannian theorem applies. The resulting
energy estimate transfers back to the original invariant mass.
A positive conformal shell excludes null and spacelike ADM vectors,
and truncating the other ends completes the endwise statement.

Section~\ref{four:sec:end} proves the end reduction;
Section~\ref{four:sec:geometry} constructs the threshold exterior and collars.
Sections~\ref{four:sec:system} and~\ref{four:sec:bounds} establish the equations,
floor exclusion and global bounds. Sections~\ref{four:sec:regularity}
and~\ref{four:sec:existence} prove regularity, existence and the final
Riemannian comparison. The shared enclosure results in
Section~\ref{rig4:sec:enclosures} supply both the minimal boundary for
that comparison and the one-sided area derivative used at equality.

The equality proof then returns to the original data.
Sections~\ref{rig4:var:section} and~\ref{rig4:adj:section} turn the
numerical inequality for nearby DEC and trapped data into a causal
stationary field. Those nearby data need not preserve outermostness
or outer area minimization. Sections~\ref{rig4:sta:section}
and~\ref{rig4:cla:section} remove twist, establish completeness and
classify the static base. Section~\ref{rig4:rec:section} recovers the
original metric and second fundamental form in regular horizon
coordinates and checks the converse, including boosted ends.

\begin{remark}[Order of constants and limits]\label{four:rem:constants}
Fix the geometry, strictification and $\epsilon>0$ first. The integer
$n$ is then sufficiently large. A bound $\Pi_n$ is polynomial in $n$,
uniformly in the outer truncation radius, homotopy parameter and solution;
fixed-$n$ regularity constants may depend arbitrarily on $n$.
Exhaust the outer radius at fixed $n$, then send $n\to\infty$ using
only the polynomial bounds. Finally send $\epsilon\downarrow0$,
remove strictification at each fixed repaired end, and send the
end-replacement radius to infinity.
\end{remark}
